% concatenated from No-Wandering.tex
\documentclass[11pt,reqno]{amsart}
\usepackage[T1]{fontenc}
\usepackage{lmodern,amsmath,amssymb,amsthm,mathtools,mathrsfs}
\usepackage[a4paper,margin=27mm]{geometry}
\usepackage{microtype,enumitem}
\usepackage[hidelinks,pdfusetitle]{hyperref}
\hypersetup{
 pdftitle={A new proof of no-wandering domain theorems without quasiconformal techniques},
 pdfsubject={Hyperbolic area proofs of no-wandering theorems without quasiconformal techniques},
 pdfkeywords={wandering domain, hyperbolic area, singular value, Poincare metric, rational map, entire function}
}
\setlength{\emergencystretch}{1.5em}
\raggedbottom
\setlist[enumerate]{leftmargin=2em,itemsep=2pt}
\newcommand{\C}{\mathbb C}
\newcommand{\D}{\mathbb D}
\newcommand{\R}{\mathbb R}
\newcommand{\Z}{\mathbb Z}
\newcommand{\sphere}{\widehat{\mathbb C}}
\newcommand{\Fat}{\mathcal F}
\newcommand{\Jul}{\mathcal J}
\newcommand{\dd}{\,\mathrm d}
\newtheorem{theorem}{Theorem}[section]
\newtheorem{proposition}[theorem]{Proposition}
\newtheorem{lemma}[theorem]{Lemma}
\newtheorem{corollary}[theorem]{Corollary}
\theoremstyle{definition}
\newtheorem{definition}[theorem]{Definition}
\newtheorem{example}[theorem]{Example}
\numberwithin{equation}{section}
\title[A new proof of no-wandering domain theorems without quasiconformal techniques]{A new proof of no-wandering domain theorems\\
without quasiconformal techniques}

% Author information: fill the braces; leave optional entries empty.
% Before submission, complete the author, address, and email fields below.
\author{Zihao Ye}   % Full name in English or Romanized form.
\address{Zihao Ye: School of Mathematical Sciences, East China Normal University, Shanghai 200241, China}  % Department, institution, city, postal code, country.
\email{52290155037@stu.ecnu.edu.cn}    % Email address.
% The PDF author metadata is filled automatically from \author above.

\date{September 20, 2026}
\subjclass[2020]{Primary 37F10; Secondary 30D05, 30F45}
\keywords{wandering domain, hyperbolic area, singular value, Poincar\'e metric, rational map, entire function}
\begin{document}
\begin{abstract}
In one complex dimension, we use hyperbolic area to prove the absence
of wandering Fatou components for rational maps of degree at least two
and for transcendental entire functions whose singular sets are compact
and have all their accumulation points in the Fatou set. As corollaries, we recover
Sullivan's theorem for rational maps and the theorem of
Eremenko--Lyubich and Goldberg--Keen for entire functions with finitely
many singular values. The proof uses neither quasiconformal deformation
nor Teichm\"uller theory.
\end{abstract}
\maketitle

\section{Introduction}

In one complex dimension, Sullivan proved that rational maps of degree
at least two on the Riemann sphere have no wandering Fatou components
\cite{Sullivan}. Eremenko--Lyubich \cite[Theorem~3]{EL} and,
independently, Goldberg--Keen \cite[Theorem~4.2]{GoldbergKeen} proved
the corresponding result for entire functions on $\C$ with finitely many
singular values.
These proofs use quasiconformal techniques.

In this paper, we develop a hyperbolic-area method that proves the
absence of wandering Fatou components for rational maps of degree at
least two and for transcendental entire functions whose singular sets
are compact and have all their accumulation points in the Fatou set.
The classical theorems above follow as corollaries, with proofs that
use no quasiconformal techniques. This gives a partial affirmative
answer to Bergweiler's Question~9 \cite[Section~4.6]{Bergweiler},
covering the rational and finite-type entire cases.

Write $\Fat(f)$ and $\Jul(f)$ for the Fatou and Julia sets. A Fatou
component is \emph{wandering} if its forward images lie in pairwise
distinct Fatou components.

\begin{definition}\label{def:singular}
For a transcendental entire function $f:\C\to\C$, let $S(f)$ be the closure in
$\C$ of its finite critical and asymptotic values.
Write $\mathcal B$ for the class with
compact $S(f)$ and $\mathcal S$ for the subclass with finite $S(f)$.
For a rational map $f:\sphere\to\sphere$, let $S(f)$ be its critical values.
In either case, $S(f)'$ denotes the set of accumulation points of $S(f)$
in the corresponding target space.
\end{definition}

\begin{theorem}[Main theorem]\label{core:singular-accumulation}
Let $f:\sphere\to\sphere$ be a rational map of degree at least two,
or let $f:\C\to\C$ be a transcendental entire function whose singular
set $S(f)$ is compact and satisfies $S(f)'\subset\Fat(f)$.
Then $f$ has no wandering Fatou components.
\end{theorem}

For compact $S(f)$, the hypothesis says equivalently that
$S(f)\cap\Fat(f)$ is compact and $S(f)\cap\Jul(f)$ is finite. It
imposes no restriction on the forward orbits of the Julia singular
values. The Fatou components containing $S(f)'$ are not assumed to be
preperiodic. The following application has infinitely many singular
values in a rotation domain.

\begin{example}\label{example:sine}
Let $\lambda=e^{2\pi i\theta}$, where $\theta$ is Diophantine irrational,
and define the entire function $h_\lambda:\C\to\C$ by
\[
 h_\lambda(z)=\lambda\frac{\sin^2z}{z},\qquad h_\lambda(0)=0.
\]
Then $h_\lambda\in\mathcal B\setminus\mathcal S$, its only finite
asymptotic value is zero, and
\begin{equation}\label{eq:sine-singular}
 S(h_\lambda)=\{0\}\cup
 \left\{\frac{4\lambda x}{1+4x^2}:
 x\in\R\setminus\{0\},\ \tan x=2x\right\},
 \qquad S(h_\lambda)'=\{0\}.
\end{equation}
The origin is a Siegel fixed point \cite{Siegel}. Hence $h_\lambda$ has no wandering
Fatou components by Theorem~\ref{core:singular-accumulation}.
\end{example}

The function-theoretic precedent is the classical example
$\sin z/z\in\mathcal B\setminus\mathcal S$
\cite[Section~4.2]{Sixsmith}. The maps in Example~\ref{example:sine}
are neither geometrically finite \cite[Proposition~5.3]{ARS}, postcritically
separated \cite[Lemma~2.6]{PS}, nor topologically hyperbolic
\cite[Definition~1.2]{BFJK}. They do not satisfy the uniformly escaping
singular-set hypothesis of \cite[Theorem~1.2]{MBRG}, and no affine
conjugate satisfies the real entire hypothesis
$g(\R)\cup S(g)\subset\R$ of \cite[Theorem~1.3]{MBRG}.

Finite singular sets have no accumulation points. We therefore recover
the classical conclusions directly.

\begin{corollary}[Sullivan]\label{core:sullivan}
Every Fatou component of a rational map $f:\sphere\to\sphere$ of
degree at least two is preperiodic.
\end{corollary}
\begin{corollary}[Eremenko--Lyubich; Goldberg--Keen]\label{core:el}
A transcendental entire function $f:\C\to\C$ with finitely many
singular values has no wandering Fatou components.
\end{corollary}
\begin{proof}[Proof of both corollaries]
Apply Theorem~\ref{core:singular-accumulation} with $S(f)'=\varnothing$.
\end{proof}

\subsection*{Proof outline}
Choose increasing finite unions $P_j$ of repelling periodic orbits,
dense in $\Jul(f)$, and put $X_j=M\setminus P_j$, where $M=\sphere$
or $\C$. Each $X_j$ has finite hyperbolic area. Let $\mu_j$ be its
normalized area form and let $\nu_j$ be the area increment obtained
by deleting $S(f)$.

Suppose a wandering component exists. The classical simple-connectivity
reductions give a simply connected tail of its orbit. In the entire
case, all Fatou components are simply connected, so we can enclose
$S(f)'$ in compact holes in finitely many of them; in the rational
case no holes are needed. Domain monotonicity and a cutoff estimate bound the
area cost of these holes uniformly in $j$, outside fixed larger
neighborhoods. The remaining singular values form a finite set, whose
area cost is bounded by its cardinality. We pass to a tail whose
components avoid these neighborhoods and singular values; on their union
$W=\bigsqcup_{n\ge0}U_n$, we have $\nu_j(W)\le B$ for all $j$.

Since each target component avoids $S(f)$, the maps
$f:U_n\to U_{n+1}$ are coverings, hence conformal isomorphisms.
Covering invariance and domain monotonicity give
\[
 \mu_j(W)\le\mu_j(W\setminus U_0)+\nu_j(W\setminus U_0).
\]
Since $\mu_j(W)<\infty$, cancellation yields $\mu_j(U_0)\le B$.
Finally, kernel convergence and monotone convergence recover the
intrinsic area of $U_0$, which is infinite because $U_0$ is simply
connected. This contradicts the uniform bound.

Section~\ref{sec:area} proves recovery of the hyperbolic metric from
finite punctures and bounds the area costs of punctures and compact
holes. Section~\ref{sec:global} uses covering invariance to establish
area cancellation along a wandering orbit and derives a finite-area
criterion. Applying the estimates of Section~\ref{sec:area} to the
singular set verifies this criterion and proves the main theorem.

\section{Hyperbolic area estimates}\label{sec:area}

\begin{definition}\label{def:area}
For a hyperbolic open set $X\subset\sphere$, let $\rho_X$ be the
complete Poincar\'e density of curvature $-4$, defined componentwise,
and let $\alpha_X$ be the normalized area measure
\[
 \dd\alpha_X=\frac2\pi\rho_X^2\,\dd x\dd y,
 \qquad \rho_\D(z)=\frac1{1-|z|^2}.
\]
\end{definition}

For hyperbolic open sets $X,Y\subset\sphere$, Schwarz--Pick and covering
invariance give
\[
 h^*\dd\alpha_Y\le\dd\alpha_X\quad(h:X\to Y\text{ holomorphic}),
 \quad h^*\dd\alpha_Y=\dd\alpha_X\quad(h\text{ a holomorphic covering}).
\]
Densities therefore increase under domain restriction. By Gauss--Bonnet,
a Riemann sphere with $N\ge3$ punctures has normalized area $N-2$; a simply
connected hyperbolic domain has infinite intrinsic area.

\begin{lemma}[Recovery from finite punctures]\label{core:kernel}
Let $D\subset\sphere$ be closed, and let $F_j\subset D$ be increasing
finite sets with $|F_1|\ge3$ and
$\overline{\bigcup_jF_j}=D$. On each component $U$ of
$\sphere\setminus D$, the densities $\rho_{\sphere\setminus F_j}$
increase locally uniformly to $\rho_U$.
\end{lemma}
\begin{proof}
For a base point in $U$, the domains $\sphere\setminus F_j$ and all
their subsequences have Carath\'eodory kernel $U$: every neighborhood of
a point of $D$ meets some $F_j$. Hejhal's kernel theorem
\cite[Theorem~1]{Hejhal} gives convergence of the normalized universal
covers, hence of the densities: at the base point the density is the
reciprocal of the absolute value of the covering map's derivative.
Monotonicity and Dini's theorem give local uniform convergence.
\end{proof}

\begin{lemma}[Area cost of punctures]\label{core:puncture}
If $X\subset\sphere$ is hyperbolic and $E\subset X$ is finite, then
\[
 \nu=\dd\alpha_{X\setminus E}-\dd\alpha_X\ge0,
 \qquad \nu(X\setminus E)\le |E|.
\]
\end{lemma}
\begin{proof}
Choose increasing finite sets $F_j\subset\sphere\setminus X$,
containing a fixed triple, with
$\overline{\bigcup_jF_j}=\sphere\setminus X$. Put
$X_j=\sphere\setminus F_j$ and
$Y_j=\sphere\setminus(F_j\cup E)$. Gauss--Bonnet gives
\[
 \int_{Y_j}(\dd\alpha_{Y_j}-\dd\alpha_{X_j})=|E|.
\]
On $X\setminus E$ both densities converge to their intrinsic limits by
Lemma~\ref{core:kernel}. Fatou's lemma for their nonnegative differences
proves the assertion, also when $X$ is disconnected.
\end{proof}

\begin{lemma}[Exterior area cost of compact holes]\label{core:holes}
Let $P_j$ be finite subsets of $\sphere$ containing a fixed triple, and
let $K\Subset\operatorname{Int}N$ be finite unions of pairwise disjoint
closed smooth Jordan disks. Suppose a fixed neighborhood of $N$ avoids
all $P_j$. For $X_j=\sphere\setminus P_j$ and $Y_j=X_j\setminus K$,
there is a constant $C<\infty$, independent of $j$, such that
\begin{equation}\label{core:holes-bound}
 \int_{X_j\setminus N}
       (\dd\alpha_{Y_j}-\dd\alpha_{X_j})\le C.
\end{equation}
\end{lemma}
\begin{proof}
The case $K=\varnothing$ is immediate. Otherwise set
$u_j=\log(\rho_{Y_j}/\rho_{X_j})\ge0$. The curvature equation gives
$\Delta u_j=4(\rho_{Y_j}^2-\rho_{X_j}^2)\ge0$.
Let $Q$ be the fixed triple, $V=\sphere\setminus Q$, and
$\varepsilon=d_V(K,V\setminus\operatorname{Int}N)>0$.
For $z\in X_j\setminus N$, domain monotonicity gives
$d_{X_j}(z,K)\ge d_V(z,K)\ge\varepsilon$.
A universal covering $\phi:\D\to X_j$ with $\phi(0)=z$ maps
$\{|w|<\tanh\varepsilon\}$ into $Y_j$. Schwarz--Pick gives
$u_j(z)\le\log\coth\varepsilon$
(cf.\ \cite[Proposition~3.4]{MBRG}).
Thus $u_j$ is locally bounded above at every point of $P_j$ and extends
to a nonnegative subharmonic function on $\sphere\setminus K$.

Let $A=\operatorname{Int}N\setminus K$. Since $A\subset Y_j$ and
$X_j\subset V$, domain monotonicity gives
\begin{equation}\label{core:holes-comparison}
 0\le u_j\le u_*:=\log\frac{\rho_A}{\rho_V}
 \quad\text{on }A.
\end{equation}

Fix a smooth conformal metric $g$ on the sphere and a smooth cutoff
$\chi\in C_c^\infty(\sphere\setminus K)$ with $0\le\chi\le1$ and
$\chi=1$ on a neighborhood of $\sphere\setminus N$.
Then $\operatorname{supp}\Delta_g\chi\Subset A$
avoids every $P_j$. The extended distributional Laplacian is positive,
so integration by parts on $\sphere\setminus K$ gives
\begin{align}\label{core:holes-cutoff}
 \int_{X_j\setminus N}(\dd\alpha_{Y_j}-\dd\alpha_{X_j})
 &\le\frac1{2\pi}\langle\Delta_g u_j,\chi\rangle
 =\frac1{2\pi}\int u_j\Delta_g\chi\,\dd A_g\notag\\
 &\le\frac1{2\pi}
 \sup_{\operatorname{supp}\Delta_g\chi}u_*
 \int|\Delta_g\chi|\,\dd A_g.
\end{align}
The last quantity is finite and independent of $j$.
\end{proof}

\section{The area criterion and the main theorem}\label{sec:global}

Let $f:M\to M$ be a rational map of degree at least two on $M=\sphere$,
or a transcendental entire function on $M=\C$.
The orbit of a Fatou component $U$ means the sequence of Fatou
components $U_n$ containing $f^n(U)$; thus $f(U_n)\subset U_{n+1}$.

\begin{definition}\label{def:marking}
Choose increasing finite unions $P_j$ of repelling periodic orbits with
\begin{equation}\label{core:marks}
 f(P_j)=P_j,\qquad |P_1|\ge3,\qquad
 \overline{\bigcup_jP_j}=\Jul(f),
\end{equation}
where closure is taken in $M$. Such sets exist by
\cite[Theorem~4]{Bergweiler}. Define
\begin{equation}\label{core:background}
 X_j=M\setminus P_j,\quad Y_j=X_j\setminus S(f),\quad
 \mu_j=\dd\alpha_{X_j},\quad
 \nu_j=\dd\alpha_{Y_j}-\mu_j\ge0\quad\text{on }Y_j.
\end{equation}
\end{definition}

Each $X_j$ is a finitely punctured sphere and has finite area.
Apply Lemma~\ref{core:kernel} with $D=\Jul(f)$ and $F_j=P_j$ in the
rational case, and with $D=\Jul(f)\cup\{\infty\}$ and
$F_j=P_j\cup\{\infty\}$ in the entire case. On every Fatou component $U$,
\begin{equation}\label{core:recovery}
 \rho_{X_j}|_U\uparrow\rho_U.
\end{equation}

\begin{lemma}[Area cancellation]\label{core:tail}
Let $U_n$ be pairwise distinct simply connected Fatou components with
$f(U_n)\subset U_{n+1}$ and $U_n\cap S(f)=\varnothing$.
Put $W=\bigsqcup_{n\ge0}U_n$. Then $f:W\to W\setminus U_0$ is a
conformal isomorphism and, for every $j$,
\begin{equation}\label{core:lossbound}
 \mu_j(U_0)\le\nu_j(W\setminus U_0).
\end{equation}
\end{lemma}
\begin{proof}
The component of $f^{-1}(U_{n+1})$ containing $U_n$ lies in $\Fat(f)$
by complete invariance, so it equals $U_n$. Since
$U_{n+1}\cap S(f)=\varnothing$, the restriction $f:U_n\to U_{n+1}$
is a surjective covering \cite[Section~3.2]{Sixsmith}.
The target is simply connected, so this covering is a conformal
isomorphism. The target components are pairwise distinct, giving the
asserted isomorphism on $W$.

Since $f(P_j)=P_j$, we have $f^{-1}(Y_j)\subset X_j$.
The map $f:f^{-1}(Y_j)\to Y_j$ is an unramified covering, so
covering invariance and domain monotonicity give
\begin{equation}\label{core:cover}
 f^*(\mu_j+\nu_j)=\dd\alpha_{f^{-1}(Y_j)}\ge\mu_j
 \quad\text{on }f^{-1}(Y_j).
\end{equation}
Since $W\subset f^{-1}(Y_j)$, integrating \eqref{core:cover} and
changing variables on $W$ gives
\[
 \mu_j(W)\le\mu_j(W\setminus U_0)+\nu_j(W\setminus U_0).
\]
Since $\mu_j(W)<\infty$, cancellation proves \eqref{core:lossbound}.
\end{proof}

\begin{proposition}[Finite-area criterion]\label{core:global}
For $f:M\to M$ as above, assume that every wandering orbit of Fatou
components has a simply connected tail. Suppose a Borel set
$Z\subset X_1$ satisfies:
\begin{enumerate}
 \item all but finitely many Fatou components meeting $M\setminus Z$
 are preperiodic;
 \item $S(f)\cap Z$ is finite and
 $\sup_j\nu_j(Z\cap Y_j)\le B<\infty$.
\end{enumerate}
Then $f$ has no wandering Fatou components.
\end{proposition}
\begin{proof}
Suppose $U$ is wandering, with component orbit $(U_n)$.
No $U_n$ is preperiodic, and the $U_n$ are pairwise distinct, so (1)
implies that only finitely many meet $M\setminus Z$.
After discarding finitely many iterates and reindexing, the $U_n$ are
simply connected, lie in $Z$, and avoid the finite set $S(f)\cap Z$.
Lemma~\ref{core:tail} gives $\mu_j(U_0)\le B$ for every $j$.
By \eqref{core:recovery} and monotone convergence, these areas tend
to the infinite intrinsic area of the simply connected domain $U_0$,
a contradiction.
\end{proof}

\begin{proof}[Proof of Theorem~\ref{core:singular-accumulation}]
Every rational wandering orbit has a simply connected tail by Baker's
covering argument \cite[Lemma~5.34, p.~90]{McMullen}.
For a transcendental entire function $f\in\mathcal B$, every Fatou
component is simply connected \cite[Proposition~3, p.~993]{EL}.

The compact set $S(f)'\Subset\Fat(f)$ meets finitely many Fatou
components. Choose finite unions of pairwise disjoint closed smooth
Jordan disks $K\Subset\operatorname{Int}N\Subset\Fat(f)$ with
$S(f)'\subset\operatorname{Int}K$. In the entire case, take images of
concentric closed disks under Riemann maps of the relevant components;
in the rational case take $K=N=\varnothing$. The set
$E_0=S(f)\setminus\operatorname{Int}K$ is finite, since otherwise
compactness would give an accumulation point of $S(f)$ outside
$\operatorname{Int}K$.

Set $Z=X_1\setminus N$. Only finitely many Fatou components meet
$M\setminus Z=P_1\cup N$, so condition (1) of
Proposition~\ref{core:global} holds. Since $S(f)\subset K\cup E_0$,
domain monotonicity gives, on
$Z\cap Y_j\subset X_j\setminus(K\cup E_0)$,
\[
 0\le\nu_j
 \le (\dd\alpha_{X_j\setminus K}-\dd\alpha_{X_j})
    +(\dd\alpha_{X_j\setminus(K\cup E_0)}
                              -\dd\alpha_{X_j\setminus K}).
\]
Since $N\Subset\Fat(f)$, a fixed neighborhood of $N$ avoids all $P_j$
and, in the entire case, infinity. Thus Lemma~\ref{core:holes} bounds
the first integral by a constant $C$ independent of $j$, using
$P_j\cup\{\infty\}$ in the entire case. Applying
Lemma~\ref{core:puncture} to $X_j\setminus K$ bounds the second by
$|E_0\cap(X_j\setminus K)|\le|E_0|$.
Hence $\nu_j(Z\cap Y_j)\le C+|E_0|$ and $S(f)\cap Z\subset E_0$,
so Proposition~\ref{core:global} applies.
\end{proof}

\subsection*{Use of artificial intelligence}
The author formulated the initial idea and guiding intuition and
established the research strategy. Working within this strategy and
under the author's direction, ChatGPT (OpenAI) identified and developed
the concrete arguments used in the proofs. The author also used ChatGPT
to assist with organizing and revising the manuscript and checking the
mathematical arguments. The author takes full responsibility for the
mathematical content, the references, and the final presentation of the
paper.

\begin{thebibliography}{99}

\bibitem{ARS}
M.~Alhamed, L.~Rempe and D.~Sixsmith, Geometrically finite transcendental entire functions,
\emph{J. Lond. Math. Soc.} (2) \textbf{106} (2022), 485--527.
\href{https://doi.org/10.1112/jlms.12516}{doi:10.1112/jlms.12516}.

\bibitem{BFJK}
K.~Bara\'nski, N.~Fagella, X.~Jarque and B.~Karpi\'nska, Fatou components and singularities of meromorphic functions,
\emph{Proc. Roy. Soc. Edinburgh Sect. A} \textbf{150} (2020), 633--654.
\href{https://doi.org/10.1017/prm.2018.142}{doi:10.1017/prm.2018.142}.

\bibitem{Bergweiler}
W.~Bergweiler, Iteration of meromorphic functions,
\emph{Bull. Amer. Math. Soc. (N.S.)} \textbf{29} (1993), no.~2, 151--188.
\href{https://doi.org/10.1090/S0273-0979-1993-00432-4}{doi:10.1090/S0273-0979-1993-00432-4}.

\bibitem{EL}
A.~E.~Eremenko and M.~Yu.~Lyubich, Dynamical properties of some classes of entire functions,
\emph{Ann. Inst. Fourier (Grenoble)} \textbf{42} (1992), no.~4, 989--1020.
\href{https://doi.org/10.5802/aif.1318}{doi:10.5802/aif.1318}.

\bibitem{GoldbergKeen}
L.~R.~Goldberg and L.~Keen, A finiteness theorem for a dynamical class of entire functions,
\emph{Ergodic Theory Dynam. Systems} \textbf{6} (1986), no.~2, 183--192.
\href{https://doi.org/10.1017/S0143385700003394}{doi:10.1017/S0143385700003394}.

\bibitem{Hejhal}
D.~A.~Hejhal, Universal covering maps for variable regions,
\emph{Math. Z.} \textbf{137} (1974), 7--20.\newline
\href{https://doi.org/10.1007/BF01213931}{doi:10.1007/BF01213931}.

\bibitem{McMullen}
C.~T.~McMullen, \emph{Riemann surfaces, dynamics and geometry},
course notes from lectures 1990--2020, Section~5.11.
\href{https://people.math.harvard.edu/~ctm/papers/home/text/class/notes/rs/course.pdf}{Author's notes}.

\bibitem{MBRG}
H.~Mihaljevi\'c-Brandt and L.~Rempe-Gillen, Absence of wandering domains for some real entire functions with bounded singular sets,
\emph{Math. Ann.} \textbf{357} (2013), no.~4, 1577--1604.
\href{https://doi.org/10.1007/s00208-013-0936-z}{doi:10.1007/s00208-013-0936-z}.

\bibitem{PS}
L.~Pardo-Sim\'on, Orbifold expansion and entire functions with bounded Fatou components,
\emph{Ergodic Theory Dynam. Systems} \textbf{42} (2022), no.~5, 1807--1846.
\href{https://doi.org/10.1017/etds.2020.147}{doi:10.1017/etds.2020.147}.

\bibitem{Siegel}
C.~L.~Siegel, Iteration of analytic functions,
\emph{Ann. of Math.} (2) \textbf{43} (1942), no.~4, 607--612.

\bibitem{Sixsmith}
D.~J.~Sixsmith, Dynamics in the Eremenko--Lyubich class,
\emph{Conform. Geom. Dyn.} \textbf{22} (2018), 185--224.
\href{https://doi.org/10.1090/ecgd/324}{doi:10.1090/ecgd/324}.

\bibitem{Sullivan}
D.~Sullivan, Quasiconformal homeomorphisms and dynamics I. Solution of the Fatou--Julia problem on wandering domains,
\emph{Ann. of Math.} (2) \textbf{122} (1985), no.~2, 401--418.
\href{https://doi.org/10.2307/1971308}{doi:10.2307/1971308}.
\end{thebibliography}

\end{document}
